\chapter{SQL Anal\'itico Avanzado: Funciones de Ventana y Agregaci\'on Multidimensional}
\label{cap:sql_analitico_window_functions}

\section{Marco Conceptual y Te\'orico (Curr\'iculo Universitario)}

Las Funciones de Ventana (\textit{Window Functions}) permiten realizar c\'alculos sobre un conjunto de filas relacionadas con la fila actual sin colapsar el n\'umero de tuplas en el resultado (a diferencia de un \texttt{GROUP BY} tradicional).

\subsection{Sintaxis Formal: Cl\'ausula \texttt{OVER()}}
\begin{lstlisting}[style=sqlstyle]
FUNCTION() OVER (
    [PARTITION BY col1, col2, ...]
    [ORDER BY col_sort ASC|DESC]
    [ROWS|RANGE frame_specification]
)
\end{lstlisting}

\subsection{Familias de Funciones Anal\'iticas}
\begin{itemize}
    \item \textbf{Ranking y Paginaci\'on:} \texttt{ROW\_NUMBER()}, \texttt{RANK()}, \texttt{DENSE\_RANK()}, \texttt{NTILE(k)}.
    \item \textbf{Desplazamiento Relativo:} \texttt{LAG(col, offset)}, \texttt{LEAD(col, offset)}.
    \item \textbf{Agregaci\'on Acumulada:} \texttt{SUM(col) OVER (...)}, \texttt{AVG(col) OVER (...)}.
\end{itemize}

---

\section{Taller de Consolidaci\'on del Cap\'itulo}

\begin{businessproblem}
El \'area financiera requiere calcular la progresi\'on mensual de ventas acumuladas por categor\'ia de producto durante todo el a\~no hist\'orico y el ranking de ventas de cada producto respecto a su categor\'ia en el mismo mes.
\end{businessproblem}

\subsection{Soluci\'on T\'ecnica en PostgreSQL (Formato O'Reilly)}

\begin{lstlisting}[style=sqlstyle, caption={Consulta anal\'itica con Window Functions y acumulados}]
WITH monthly_product_revenue AS (
    SELECT
        p.category_id,
        c.category_name,
        p.product_name,
        DATE_TRUNC('month', o.order_date)::DATE AS sales_month,
        SUM(od.unit_price * od.quantity * (1 - od.discount)) AS total_revenue
    FROM orders AS o
    INNER JOIN order_details AS od
        ON o.order_id = od.order_id
    INNER JOIN products AS p
        ON od.product_id = p.product_id
    INNER JOIN categories AS c
        ON p.category_id = c.category_id
    GROUP BY
        p.category_id,
        c.category_name,
        p.product_name,
        DATE_TRUNC('month', o.order_date)::DATE
)
SELECT
    category_name,
    product_name,
    sales_month,
    total_revenue::NUMERIC(10, 2) AS monthly_sales,
    -- Acumulado m\'ovil anual por producto
    SUM(total_revenue) OVER (
        PARTITION BY category_name, product_name
        ORDER BY sales_month
    )::NUMERIC(12, 2) AS running_total_ytd,
    -- Ranking mensual dentro de la categor\'ia
    DENSE_RANK() OVER (
        PARTITION BY category_name, sales_month
        ORDER BY total_revenue DESC
    ) AS category_monthly_rank
FROM monthly_product_revenue
ORDER BY
    category_name ASC,
    sales_month ASC,
    category_monthly_rank ASC;
\end{lstlisting}

\begin{engtip}
\begin{itemize}
    \item \textbf{\texttt{DENSE\_RANK()} vs \texttt{RANK()}:} \texttt{DENSE\_RANK()} asigna rangos consecutivos sin saltos num\'ericos en caso de empates de montos ($1, 2, 2, 3...$), ideal para tablas de clasificaci\'on.
\end{itemize}
\end{engtip}
